%-------------------------------------------------------------
% One-page resume for internship applications
% Syed Emad Uddin Shubha
% Compile with pdflatex
%-------------------------------------------------------------

\documentclass[letterpaper,10pt]{article}

\usepackage[margin=0.5in,top=0.4in,bottom=0.4in]{geometry}
\usepackage[T1]{fontenc}
\usepackage{charter}
\usepackage{titlesec}
\usepackage{enumitem}
\usepackage[dvipsnames]{xcolor}
\usepackage[hidelinks]{hyperref}
\usepackage{tabularx}
\usepackage{amsmath}
\usepackage[english]{babel}
\input{glyphtounicode}
\pdfgentounicode=1

\hypersetup{pdftitle={Resume -- Syed Emad Uddin Shubha}, pdfauthor={Syed Emad Uddin Shubha}}

\pagestyle{empty}
\raggedright
\setlength{\parindent}{0pt}
\urlstyle{same}

\titleformat{\section}{\large\scshape\bfseries}{}{0em}{}[\vspace{-3pt}\titlerule]
\titlespacing*{\section}{0pt}{6pt}{3pt}

\newcommand{\me}{\textbf{S.\,E.\,U.~Shubha}}
\newcommand{\cvlink}[2]{\href{#1}{\underline{#2}}}
\newcommand{\entry}[4]{%
  \begin{tabularx}{\textwidth}{@{}X r@{}}
    \textbf{#1} & #2 \\
    \textit{#3} & \textit{#4}
  \end{tabularx}\par}

\setlist[itemize]{leftmargin=1.1em, itemsep=0pt, topsep=1pt, parsep=0pt}
\interlinepenalty=10000

\begin{document}

%---------- Header ----------
\begin{center}
  {\LARGE\scshape\bfseries Syed Emad Uddin Shubha}\\[3pt]
  \small
  \cvlink{mailto:syed.shubha@ieee.org}{syed.shubha@ieee.org} \;$|$\;
  \cvlink{https://syedshubha.com/}{syedshubha.com} \;$|$\;
  \cvlink{https://www.linkedin.com/in/syedshubha/}{in/syedshubha} \;$|$\;
  \cvlink{https://github.com/syedshubha}{github.com/syedshubha} \;$|$\;
  \cvlink{https://scholar.google.com/citations?user=8EK8VJMAAAAJ}{Google Scholar}   \\[1pt] %\cvlink{https://orcid.org/0009-0005-1435-5958}{ORCID}  \\[1pt]

  Citizenship: Bangladesh \;$|$\; US Visa Status: F-1 Student
\end{center}
\vspace{-4pt}

%---------- Summary ----------
\section{Summary}
Ph.D. student in Computer Science working on \textbf{quantum hardware security} and \textbf{quantum communication}. Builds pulse-level and circuit-level simulation and verification tools for superconducting quantum processors. First-author papers at ACM QSec 2025 and IEEE QCE 2025 (two papers). First-author crosstalk detection paper accepted for \textbf{IEEE HOST 2027} (to appear). Four journal articles (\cvlink{https://syedshubha.com/\#publications}{full publication list}). \textbf{Quantum Rising Star Award} (Rice University, 2026).

%---------- Education ----------
\section{Education}
\entry{Louisiana State University}{January 2025 -- Present}
      {Ph.D. in Computer Science \;$\cdot$\; Advisor: Dr.~Tasnuva Farheen}{GPA: 4.0/4.0}
\vspace{-2pt}
\begin{itemize}
  \item Coursework: Quantum Technology (A+), Quantum Algorithms (A+), Machine Learning (A+)
\end{itemize}
\entry{Bangladesh University of Engineering and Technology (BUET)}{February 2017 -- May 2022}
      {B.Sc. in Electrical and Electronic Engineering}{Dhaka, Bangladesh}

%---------- Research Experience ----------
\section{Research Experience}
\entry{Graduate Research Assistant, Louisiana State University}{January 2025 -- Present}
      {Quantum hardware security and quantum communication (directions I proposed)}{Baton Rouge, LA}
\vspace{-2pt}
\begin{itemize}
  \item Modeled crosstalk between co-located users on shared superconducting hardware with pulse-level simulation, and built a pulse-to-circuit framework that identified a new asymmetric attack mechanism (IEEE QCE '25; ACM QSec '25).
  \item Designed a detector that separates malicious crosstalk from benign calibration drift using Pauli-transfer-matrix signatures, built on an exact structural zero in the drift signature (IEEE HOST '27).
  \item Compared simulated operations under isolated and concurrent execution to study interference on shared quantum hardware.
  \item Designed an adaptive purification protocol for networked superdense coding (IEEE QCE '25) and studied how noise affects correlations between qubits.
  \item Conceived and wrote the research aims and methodology of an awarded NSF proposal (\$194,596) on adaptive error correction and purification for quantum networks, and now carry out its research; hardware-security work funded by a Louisiana Board of Regents award (\$73,926).
\end{itemize}
\vspace{2pt}
\entry{Research Assistant, North South University}{June 2022 -- May 2024}
      {Supervisor: Dr.~Mahdy Rahman Chowdhury}{Dhaka, Bangladesh}
\vspace{-2pt}
\begin{itemize}
  \item Pioneered Bangladesh's first research initiative in quantum computing and communication; helped build the NSU Center of Quantum Computing.
  \item Proposed a stabilizer-formalism error-correction protocol that improves Bell-pair fidelity under Pauli noise (first author, \textit{Heliyon}).
  \item Led student research projects (surface-code decoders, QKD-based security) that produced 3 papers; 5 former mentees are now Ph.D. students in quantum information.
\end{itemize}

%---------- Projects ----------
\section{Projects}
\begin{itemize}
  \item \textbf{Superconducting Qubit Simulation and Hardware Modeling} (\cvlink{https://github.com/syedshubha/EmulateQPU-4Q}{GitHub}). Built a Python pulse simulator for four qubits and three tunable couplers. Compared capacitances and charging energies across three Gmsh meshes with Palace.
  \item \textbf{W-State Preparation with Qiskit Pulse} (\cvlink{https://github.com/syedshubha/Phy4750Sh}{GitHub}): designed gate and measurement pulses from \textit{ibm\_sherbrooke} calibration data; Rabi $\pi$-pulse calibration and echoed cross-resonance to suppress crosstalk.
  \item \textbf{\cvlink{https://syedshubha.com/qsvt/}{Interactive QSVT Laboratory}}. Built a browser tool for exploring block encoding and quantum signal processing.
  \item \textbf{Neural Shadows} (\cvlink{https://github.com/syedshubha/ShadowTomography}{GitHub}): ML classification of Werner-state entanglement from sparse classical shadows; Monte Carlo robustness analysis of SVM and FNN models under finite-shot noise.
\end{itemize}

%---------- Skills ----------
\section{Skills}
\begin{tabularx}{\textwidth}{@{}l X@{}}
  \textbf{Languages \& tools:} & Python, C, C++, Rust, MATLAB, OpenQASM, Git, \LaTeX \\
  \textbf{Quantum tools:} & Qiskit (incl.\ Qiskit Pulse), QuTiP, Amazon Braket \\
  \textbf{Methods:} & Pulse-level Hamiltonian simulation, process tomography / Pauli transfer matrices, Choi-matrix \mbox{analysis}, classical shadows, stabilizer codes, QKD protocol analysis. \\
\end{tabularx}

%---------- Honors & Service ----------
\section{Honors, Service, and Teaching}
\begin{itemize}
  \item Quantum Rising Star Award, Quantum Student Success Summit, Rice University (2026); NSF Student Travel Grant, IEEE QCE (2025); Phi Kappa Phi (2026).
  \item Reviewer, \textit{IEEE Transactions on Quantum Engineering} (2026). Sub-reviewer, DATE 2027 and IEEE DCAS 2026. Graduate TA, CSC 4501 Computer Networks, LSU (2025--present); TA, Quantum Computing II, BRAC University (2023--2024).
\end{itemize}

\end{document}